\documentclass[letterpaper,12pt]{article}
\usepackage{journal-template}


%------- page header/footer formatting -----------
%          Do not change
\usepackage{fancyhdr}
\pagestyle{fancy}
\lhead{}
\chead{}
\rhead{}
\lfoot{}
\cfoot{}
\rfoot{\thepage}
\renewcommand{\headrulewidth}{0pt}



%------- Title page   -----------
%        Modify as needed
\title{Type the Title of your Paper, Capitalize First Letter}
\author[1]{First Author}
\author[2]{Second Author}
\author[1]{Third Author}
\affil[1]{First affiliation, Address, City and Postcode, country}
\affil[2]{Second affiliation, Address, City and Postcode, Country}
\date{}                     %% if you don't need date to appear
\setcounter{Maxaffil}{0}
\renewcommand\Affilfont{\itshape\small}




\begin{document}
\maketitle %Automatically generated the title page


%------- Abstract -----------
%        Modify as needed
\begin{abstract}

\vspace{0.2cm}
\noindent
{\it Maximum 250 words. Times New Roman 12 pt. font.}\\

\noindent
In scientific journals, the Abstract is meant to tell a potential reader whether or not the article contains information worth reading in full. Therefore, the Abstract presents a comprehensive factual summary using a minimum number of words. This summary should include a short (i.e. 1 - 2 sentences!) description of: 
\begin{enumerate}
\item	what was done, 
\item	how it was done (experimental technique used with the relevant experimental conditions), 
\item	what were the final results (including the error estimate),
\item	what are the main conclusions (e.g., if the results are
consistent with a particular model, literature data, etc.). 
\end{enumerate}
It should be kept in mind that since the Abstract is read together with the title of the report any information given in the title need not be repeated in the Abstract. 
  
\end{abstract}




%-------------------------------------------------------------
\section{Introduction}
{\it Maximum 500 words. Times New Roman 12 pt. font.}\\
        
This section outlines basic goal of the experiments and the relevant theory used for interpreting/analyzing the results should appear first in the main body of the report. The purpose of the Introduction is to present the background information, basic concepts and equations required by the reader to understand i) the purpose of the experiment and ii) the methods used in treating the data. 

It serves little purpose, and is not considered good practice, to
rewrite or copy large portions from notes, the lab manual or
texts. Furthermore, derivations of equations given elsewhere should
not be reproduced. If the derivation is of particular interest
reference can be made to the appropriate source or it can be placed in
an appendix. However, a brief presentation of basic theory, with
relevant equations, in the words and style of the author is
essential.\\

Your work will likely build on, rely on, or otherwise make use of the results or work of others.  In all such cases it is essential that you properly reference the appropriate
  original source.  {\bf References should be numbered consecutively in
  order of first appearance throughout the text.} Software is available
  to the McGill community to make this very easy (EndNote for
  Window/Word, BibTeX for LaTeX) which is available to all McGill
  students.  References to Web URL's are to be discouraged but are
  sometimes necessary.  This is an example of a reference to a peer-reviewed journal article~\cite{paper1}, a book~\cite{book1}, online material~\cite{link1}, and a technical document such as a Laboratory manual~\cite{LabManualX}.  For papers with more than five authors include only the first author's name followed by ’et al.’.

See additional information on this and other sections in the lecture
notes.





%-------------------------------------------------------------
\section{Methods}
{\it Maximum 500 words. Maximum 2 Figures. Times New Roman 12
  pt. font.}\\


A description of the equipment and procedures used to make the
measurements are given in this section. When special apparatus is
used, a description of the operating principles together with a
clearly labelled schematic diagram is desirable. 

Experimental procedure should include a concise outline of the
        actual procedure, which should differ somewhat from that
        described in the notes or elsewhere. This section should NOT
        be a detailed transcription of the manual or take the form of
        a ``recipe''; e.g. ``and then we pressed 'setup 2' on the
        oscilloscope, and then moved the cursor to measure the
        position of...''. Instead, it should be a very concise outline
        of the main points of the procedure sufficient to convey a
        clear overview of the experimental procedure to the
        reader. This outline should consist of concise factual
        statements, e.g., ``The heat capacity of the calorimeter and
        contents were determined by measuring the temperature change
        that resulted from the addition of a known quantity of heat''.  



%-------------------------------------------------------------
\section{Results}
{\it Maximum 500 words.  Use the \underline{minimum number} of Figures and Tables
  necessary to fully present your results. Times New Roman 12 
  pt. font.}\\


This section presents the experimental data and the results of any
analyses including the estimation of experimental uncertainties.  {\bf
  Your
results should be presented in a logical sequence.  Do not jump around
in an inconsistent, hard to follow manner.} Any assumptions made in the
treatment of data should be clearly specified.  The reader should be
informed, using explanatory sentences how the raw data obtained in the
laboratory were transformed into the desired results; i.e. your graphs
and tables should appear in this section, but this section is not ONLY
your graphs and tables.


%-------------------------------------------------------------
\section{Discussion}
{\it Maximum 500 words.  Use the \underline{minimum number} of Figures and Tables
  necessary to fully discuss your results. Times New Roman 12 
  pt. font.}\\

The Discussion evaluates whether or not the experiment objectives have
been met (i.e. was the experiment a success) and what the results
actually mean. This portion of the report allows some freedom, but the
purpose is to explain the significance/meaning of the results to the
reader; i.e. to ``interpret'' the results.

In the discussion section, important features of the results should be
highlighted. Comparison should be made with accepted values if
relevant and available.  This comparison gives an estimate of the
accuracy of the experiment and provides an opportunity to assess
different methods (i.e. those you used and those others have used) and
to comment on your error analysis. {\bf Do the results obtained in this
experiment agree - within quoted uncertainty - of accepted values? If
there is a discrepancy, by how many standard deviations?  Why?}

What (if any) is the theoretical significance of the results? The
Discussion should include a clear, concise presentation of how the
results demonstrate key concepts or impact on questions of theoretical
importance.

{\bf Any direct questions asked of you in the lab manual should also be
  addressed in this section.}

%-------------------------------------------------------------
\section{Conclusions}
{\it Maximum 300 words. Times New Roman 12 
  pt. font.}\\

In the Conclusion section, state the most important outcome of your
work. Do not simply summarize the points already made in the body,
instead, interpret your findings at a higher level of
abstraction. Show whether, or to what extent, you have succeeded in
addressing the need stated in the Introduction. At the same time, do
not focus on yourself (for example, by restating everything you
did). Rather, show what your findings mean to readers. Make the
Conclusion interesting and memorable for them.

At the end of your Conclusion, consider including perspectives, that
is, an idea of what could or should still be done in relation to the
issue addressed in the paper. If you include perspectives, clarify
whether you are referring to firm plans for yourself and your
colleagues (``In the coming months, we will . . . ``) or to an
invitation to readers (``One remaining question is . . . ``).



%-------------------------------------------------------------
\bibliographystyle{IEEEtran}
%\bibliographystyle{unsrt}
\bibliography{MyBibliography}




%-------------------------------------------------------------
\section*{Author Contributions Statement}
All manuscripts must include an Author Contributions Statement
detailing the contribution of each author. Author names must be listed
as initials. For example, ``A.B and C.D wrote the main manuscript text
and E.F prepared figures 1-3. All authors reviewed the manuscript.''



%-------------------------------------------------------------
\appendix
\renewcommand{\thesection}{Appendix \Alph{section}. }
\renewcommand{\thesubsection}{\Alph{section}.\arabic{subsection}. }



\section{An example appendix}

Authors including an appendix section should do so after References
section.

The relevant sample calculations for the data tables and figures should appear in an appendix


\subsection{Example of a sub-heading within an appendix}


\end{document}
